% Options for packages loaded elsewhere
\PassOptionsToPackage{unicode}{hyperref}
\PassOptionsToPackage{hyphens}{url}
\documentclass[
  ignorenonframetext,
]{beamer}
\newif\ifbibliography
\usepackage{pgfpages}
\setbeamertemplate{caption}[numbered]
\setbeamertemplate{caption label separator}{: }
\setbeamercolor{caption name}{fg=normal text.fg}
\beamertemplatenavigationsymbolsempty
% remove section numbering
\setbeamertemplate{part page}{
  \centering
  \begin{beamercolorbox}[sep=16pt,center]{part title}
    \usebeamerfont{part title}\insertpart\par
  \end{beamercolorbox}
}
\setbeamertemplate{section page}{
  \centering
  \begin{beamercolorbox}[sep=12pt,center]{section title}
    \usebeamerfont{section title}\insertsection\par
  \end{beamercolorbox}
}
\setbeamertemplate{subsection page}{
  \centering
  \begin{beamercolorbox}[sep=8pt,center]{subsection title}
    \usebeamerfont{subsection title}\insertsubsection\par
  \end{beamercolorbox}
}
% Prevent slide breaks in the middle of a paragraph
\widowpenalties 1 10000
\raggedbottom
\AtBeginPart{
  \frame{\partpage}
}
\AtBeginSection{
  \ifbibliography
  \else
    \frame{\sectionpage}
  \fi
}
\AtBeginSubsection{
  \frame{\subsectionpage}
}
\usepackage{iftex}
\ifPDFTeX
  \usepackage[T1]{fontenc}
  \usepackage[utf8]{inputenc}
  \usepackage{textcomp} % provide euro and other symbols
\else % if luatex or xetex
  \usepackage{unicode-math} % this also loads fontspec
  \defaultfontfeatures{Scale=MatchLowercase}
  \defaultfontfeatures[\rmfamily]{Ligatures=TeX,Scale=1}
\fi
\usepackage{lmodern}
\ifPDFTeX\else
  % xetex/luatex font selection
\fi
% Use upquote if available, for straight quotes in verbatim environments
\IfFileExists{upquote.sty}{\usepackage{upquote}}{}
\IfFileExists{microtype.sty}{% use microtype if available
  \usepackage[]{microtype}
  \UseMicrotypeSet[protrusion]{basicmath} % disable protrusion for tt fonts
}{}
\makeatletter
\@ifundefined{KOMAClassName}{% if non-KOMA class
  \IfFileExists{parskip.sty}{%
    \usepackage{parskip}
  }{% else
    \setlength{\parindent}{0pt}
    \setlength{\parskip}{6pt plus 2pt minus 1pt}}
}{% if KOMA class
  \KOMAoptions{parskip=half}}
\makeatother
\setlength{\emergencystretch}{3em} % prevent overfull lines
\providecommand{\tightlist}{%
  \setlength{\itemsep}{0pt}\setlength{\parskip}{0pt}}
% Auto-generated by infrastructure/rendering/_pdf_latex_helpers.py::extract_math_font_preamble.
% Minimal header passed to Pandoc via -H so Beamer slides pick up the same math font
% as the combined-PDF path. Adjust the manuscript's preamble.md to change the font globally.
\usepackage{unicode-math}
\setmathfont{latinmodern-math.otf}
\usepackage{bookmark}
\IfFileExists{xurl.sty}{\usepackage{xurl}}{} % add URL line breaks if available
\urlstyle{same}
\hypersetup{
  hidelinks,
  pdfcreator={LaTeX via pandoc}}

\author{\texorpdfstring{}{}}
\date{}

\begin{document}

\section{Magnitude and Magnitude Homology: Effective Role Count, Lawvere
Similarity Spaces, and Language as Enriched
Category}\label{sec:magnitude-homology}

\begin{frame}[fragile,allowframebreaks]{Magnitude and Magnitude
Homology: Effective Role Count, Lawvere Similarity Spaces, and Language
as Enriched Category}
\textbf{Where we are in the argument.} \autoref{sec:enriched-categories}
introduced \([0,1]\)-enrichment so that the case category carries a full
distributional-proximity matrix. This chapter extracts the principal
quantitative invariant of that matrix --- Leinster's \emph{magnitude}
\(|\mathcal{C}| = \sum_{i,j} (Z^{-1})_{ij}\) --- which gives a
single-number answer to ``how many effectively distinct case roles does
the language use?'' (\(|\mathcal{C}| \approx 2.50\) for the standard
eight-case system), and which will serve as the N400 magnitude-change
proxy in \autoref{sec:diagrammatic-cognition} and as a
cognitive-security invariant in \autoref{sec:cognitive-security}.

A central mathematical invariant unique to enriched categories is their
\textbf{magnitude}---a numerical quantity capturing the ``effective
size'' of the category by discounting distributional overlap between
objects.

For an enriched category with \(n\) objects, let \(Z\) be the
\(n \times n\) similarity matrix where \(Z_{ij} = \mathcal{C}(i, j)\).
The categorical magnitude is:

\begin{equation}
|\mathcal{C}| = \sum_{i,j} (Z^{-1})_{ij}
\label{eq:eq-5-3}
\end{equation}

assuming \(Z\) is invertible. This magnitude metric connects to
information theory:

\begin{itemize}
\tightlist
\item
  For a \emph{discrete} category (no non-trivial relationships),
  \(|\mathcal{C}| = n\) (the number of objects)
\item
  For a highly connected category, \(|\mathcal{C}| < n\) (objects are
  ``redundant'')
\item
  Magnitude connects to the diversity measures studied in ecology
  (species diversity), graph theory (effective graph resistance), and
  information geometry (Fisher information)
\end{itemize}

\textbf{Worked example.} Consider the minimal 3-case category with
objects \(\{S, A, P\}\) and hom-values \(\mathcal{C}(S,A) = 0.85\) (S
and A share agentive contexts), \(\mathcal{C}(A,P) = 0.70\) (transitive
co-occurrence), \(\mathcal{C}(S,P) = 0.40\) (weak S--P overlap). The
similarity matrix \(Z\) is:

\begin{equation}
Z = \begin{pmatrix} 1.00 & 0.85 & 0.40 \\ 0.85 & 1.00 & 0.70 \\ 0.40 & 0.70 & 1.00 \end{pmatrix}, \quad Z^{-1} \approx \begin{pmatrix} 4.93 & -5.51 & 1.88 \\ -5.51 & 8.12 & -3.48 \\ 1.88 & -3.48 & 2.68 \end{pmatrix}
\label{eq:eq-5-4}
\end{equation}

The magnitude is
\(|\mathcal{C}| = \sum_{i,j} (Z^{-1})_{ij} \approx 4.93 - 5.51 + 1.88 - 5.51 + 8.12 - 3.48 + 1.88 - 3.48 + 2.68 \approx 1.52\)---substantially
less than the cardinality 3 of the role inventory. The deficit
\(3 - 1.52 = 1.48\) is substantial---approximately 49\% of the
cardinality---reflecting that S and A share dense agentive contexts
(\(w = 0.85\)) and A and P share transitive co-occurrence contexts
(\(w = 0.70\)), making the proto-role system considerably redundant. By
contrast, an accusative alignment---which merges S and A into a single
NOM role---would yield a 2-object category with magnitude exactly 2.0,
and the deficit \(3 - 2.0 = 1.0\) quantifies the information lost by
neutralization.

\textbf{Scaling to full categories.} Our \texttt{EnrichedCategory}
implementation computes magnitude for any case category. For the
standard 8-case English category with empirically calibrated
distributional proximity values, the magnitude is approximately
2.50---the deficit \(8 - 2.50 = 5.50\) reflects substantial
distributional overlap: NOM and ACC share transitive contexts
(\(w = 0.85\)), DAT and ACC overlap in double-object constructions
(\(w = 0.55\)), GEN and NOM co-occur in possessive constructions
(\(w = 0.60\)), and even peripheral cases such as VOC overlap with NOM
(\(w = 0.70\)). Only 2.50 of the 8 case roles encode genuinely
independent relational distinctions. This magnitude differential
provides a quantitative formalization of Silverstein's
{[}-@silverstein1976hierarchy{]} case hierarchy: languages with more
alignment-based neutralization (lower magnitude) have less relational
discriminability, while richer case inventories (higher magnitude) make
finer-grained relational distinctions.

Bradley {[}-@bradley2021entropy{]} established a link connecting
categorical magnitude to classical information entropy via topological
operad derivations. Her result proves that Shannon entropy acts as the
unique algebraic derivation of a specific topological operad---a
categorical structure governing the composition of enriched categories.
This supplies theoretical justification for magnitude as a measurable
geometric invariant quantifying linguistic complexity: the magnitude of
any case category quantifies how much irreducible ``information'' that
case system encodes regarding relational meaning.

Leinster and Shulman {[}-@leinster2021magnitude{]} further develop
\textbf{magnitude homology}, which categorifies magnitude from a scalar
invariant to a graded homological invariant---detecting not just the
``effective number'' of objects but the higher-dimensional ``holes'' in
the distributional landscape. For case categories, magnitude homology
can distinguish between two systems with the same magnitude but
different topological structure: a category where NOM--ACC--DAT form a
tight cluster and all other cases are isolated looks identical in
magnitude to one where the clustering is evenly distributed, but their
magnitude homology groups differ, revealing that the former has a
1-dimensional ``hole'' (a missing transitive link) that the latter
fills. This finer invariant provides a richer classification of case
systems than magnitude alone. Bradley and Vigneaux
{[}-@bradleyvigneaux2025magnitude{]} realize this programme on natural
language by building categories of texts enriched from language-model
next-token probabilities and computing magnitude and magnitude homology
for associated metric spaces of texts---a concrete large-scale
application of Leinster--Shulman theory beyond finite toy examples.

However, that LM-enriched construction exposes a vulnerability relevant
to our cognitive synthesis: if magnitude homology computations are
ported into non-classical environments, such as lambeq Gen II's
Parameterized Quantum Circuits (\autoref{sec:quantum-active-inference}),
the inherent environmental quantum noise (decoherence) acts as a
non-trivial topological perturbation. Unless error-correction syndromes
explicitly preserve the sequence of homology groups, the invariant may
technically fail to commute. Thus, comparing the 1-dimensional ``holes''
of classical case-frames against their quantum algorithmic analogues
must be approached with mathematical caution, respecting the shear
forces introduced by measurement non-commutativity.
\end{frame}

\begin{frame}[allowframebreaks]{Language as Enriched Category:
Transformer Attention Weights Are Context-Dependent Hom-Values}
\protect\phantomsection\label{language-as-enriched-category-transformer-attention-weights-are-context-dependent-hom-values}
Bradley's {[}-@bradley2024ipam; -@bradley2025tea{]} broader program
treats natural language itself as an enriched category, where:

\begin{itemize}
\tightlist
\item
  \textbf{Objects} are expressions (words, phrases, sentences)
\item
  \textbf{Hom-values} encode distributional co-occurrence probabilities
\item
  \textbf{Composition} models transitivity of distributional relatedness
\end{itemize}

This ``language as enriched category'' perspective has profound
implications for case theory:

\begin{enumerate}
\item
  \textbf{Case roles emerge from distributional structure}: Rather than
  being imposed a priori, case distinctions arise from clusters of high
  hom-values in the enriched language category. Nouns that frequently
  appear in agent contexts cluster together, forming the ``nominative''
  region of the category.
\item
  \textbf{Alignment types correspond to enriched structure}: Different
  languages partition the enriched category differently, and these
  partitions correspond to the alignment types (accusative, ergative,
  etc.) discussed in \autoref{sec:case-systems}.
\item
  \textbf{Language models are enriched functors}: A neural language
  model (such as a transformer) can be viewed as an enriched functor
  from the syntactic category to the semantic category, mapping
  type-logical derivations to distributional meaning representations
  while preserving the enriched structure.
\end{enumerate}

The deep connection to modern distributional semantics is this: static
embeddings operationalize hom-values as cosine similarity in a learned
space, while contextualized transformers {[}@devlin2019bert;
@vaswani2017attention{]} compute \emph{dynamic} hom-values from
sentential context---a move from a fixed enriched category to a
parameterized one. The attention-as-enriched-cup analogy (developed in
\autoref{sec:categorical-semantics}) carries the same intuition here:
layer-wise weights grade how strongly tokens couple, alongside Bradley
et al.'s {[}-@fritz2021enriched{]} probabilistic reading of hom-objects.
\end{frame}

\begin{frame}[fragile,allowframebreaks]{Lawvere's Insight: Case
Categories Are Similarity Spaces}
\protect\phantomsection\label{lawveres-insight-case-categories-are-similarity-spaces}
The \([0,1]\)-enrichment connects to a deep tradition in categorical
algebra. Lawvere showed that metric spaces are categories enriched over
\(([0, \infty], +, 0)\): the hom-value is the distance between points,
the identity axiom says \(d(x, x) = 0\), and the composition inequality
is the triangle inequality \(d(x, z) \leq d(x, y) + d(y, z)\). Our
\([0,1]\)-enrichment is the multiplicative analogue: hom-values are
\emph{similarities} rather than distances, and the composition
inequality is sub-multiplicative rather than sub-additive. The
inequality \emph{direction reverses} between the two settings because
the monoidal structure reverses: in the additive metric setting we want
small composites (triangle inequality is an upper bound on distance),
whereas for multiplicative similarity we want large composites (so the
natural inequality is the lower bound
\(\mathcal{C}(A,C) \ge \mathcal{C}(A,B)\cdot\mathcal{C}(B,C)\),
equivalently the upper-bound form used in the \autoref{sec:notation}
notation appendix).

\textbf{When is magnitude defined?} Leinster's magnitude requires the
similarity matrix \(Z\) to be invertible. For the standard eight-case
category (\autoref{tbl:eight-cases}) \(Z\) has condition number
\(\approx 9.5\) and magnitude is well-defined (\(|C| \approx 2.50\),
deficit \(5.50\)). For degenerate categories where roles are
distributional clones (e.g.~\(Z\) has repeated rows), \(Z\) becomes
singular; the implementation in \texttt{src/enriched\_cat/enriched.py}
falls back to a Moore--Penrose pseudo-inverse with a logged warning,
which yields a best-least-squares approximate magnitude but not the
exact Leinster magnitude (which is simply undefined in that case).

\textbf{Distributional-semantics embedding into \([0,1]\).} A
word-embedding model instantiates the enriched category via the explicit
map
\(\mathcal{C}(w, w') = \bigl(\cos(v_w, v_{w'}) + 1\bigr) / 2 \in [0,1]\),
where \(v_w, v_{w'}\) are the model's vector representations and
\(\cos\) is cosine similarity. This is the concrete choice that allows a
pretrained transformer or word2vec model to be read as a
\([0,1]\)-enriched functor.

This Lawvere-style perspective unifies our case categories with the
geometry of distributional semantics: case roles are points in a
``similarity space,'' and morphisms between them are paths weighted by
distributional proximity. The magnitude of this space then quantifies
the ``effective dimensionality'' of the case system---how many
independent relational distinctions the language makes.
\end{frame}

\end{document}
